\paragraph{Ground occupation and energy constraints under refinement.}
The original zero-chemical-potential ground satisfies a useful uniform
constraint before any continuum particle identification. With normalized
Brillouin measure, let $\rho(k)$ be its one-body covariance and $P_\pm(k)$
the spectral projectors of the original Weyl hopping. The free filled sea
has onsite covariance $I/2$ and double occupation $1/4$. Its finite-box
restrictions give the variational bound
\[
 0\le e_{\rm kin}-e_{\rm free}+Ud=e_U-e_{\rm free}\le U/4,
 \qquad d=\omega(n_\uparrow n_\downarrow).
\]
Here $e_{\rm free}=-c\int\sqrt{\sum_j\sin^2 k_j}$; no free reference is
substituted for the interacting state. Set
\[
 w(k)=\operatorname{tr}(P_+\rho)+\operatorname{tr}(P_-(I-\rho)),
 \qquad D=\int w(k).
\]
Then $0\le w\le2$ and the kinetic excess is
$c\int\sqrt{\sum_j\sin^2 k_j}\,w(k)$. The region with
$\sqrt{\sum_j\sin^2 k_j}\le\delta$ has normalized volume at most
$\delta^3$ for $0<\delta\le1$, since
$2\arcsin\delta/\pi\le\delta$. Splitting this integral proves
\[
 D\le\min\{2,2\delta^3+U/(4c\delta)\},\qquad |n-1|\le D.
\]
For $U\le24c$, choosing $\delta=(U/(24c))^{1/4}$ gives
$D\le8(U/(24c))^{3/4}$. Thus the density tends to one per site,
and the per-spin density $r$ tends to $1/2$, when $U/c\to0$.
At $v=1/8$, $U=1$, $a=1/65536$, $c=v/a$ and $\delta=1/16$,
the exact bound is $1023/1024\le n\le1$.

The same defect controls every fixed microscopic local word.
The free-sea annihilation parts $c(P_+f)$ and $c^\dagger(P_-f)$ of a
unit site-basis field have interacting GNS norm at most $\sqrt D$.
Moving the first annihilation part to the right by CAR gives the free
Wick recursion, with two boundary errors at most $\sqrt D$. Hence for
an even word of length $m$,
\[
 |\omega(A_1\cdots A_m)-\omega_{\rm free}(A_1\cdots A_m)|
 \le C_m\sqrt D,\qquad C_0=0,\quad C_m=2+(m-1)C_{m-2}.
\]
In particular $C_4=8$ and $d\to1/4$. This controls fixed microscopic
local observables; shrinking momentum pockets can still carry a full
normalized continuum packet, so no Gaussian continuum conclusion follows.

For cell-averaged physical $H^1$ packets and $c=v/a$, bounded
$U(a)\to U_0$ fixes the first spectral moment to
$\langle f,h_{\rm W}f\rangle+(U_0/2)\|f\|^2$ in every weak spectral
limit. The uniform second-moment bound supplies the required first-moment
tail control. The scalar term is retained without a chemical-potential
shift. Second moments need not pass: the positive measures
$(1-n^{-2})\delta_0+n^{-2}\delta_n$ converge weakly to $\delta_0$
while their variances tend to one. Thus even a nonzero memory Gram need
not survive the limiting response. These constraints supply ground and
energy inputs, not relativistic rest masses or an Einstein balance.
The proof and exact controls are in
\texttt{model/continuum\_ground\_constraints}.
